\renewcommand{\Anno}{2011}

\renewcommand{\NumeroProblema}{4}

\renewcommand{\Testo}{
\begin{enumerate}
	\item  
	Un gruppo di 8 amici ha acquistato dei biglietti di teatro,
	numerati da 1 a 8, che occupano un'intera fila.  Giunti in sala,
	si siedono nella fila prenotata ma in un ordine casuale.

	Determinare la probabilit\`{a} che ognuno sieda su una sedia il
	cui numero differisce da quello del suo biglietto per al pi\`{u}
	uno.

	\item Dopo il teatro, gli 8 amici hanno prenotato in un ristorante
	per la cena.  Il proprietario ha riservato per loro un tavolo
	rotondo con 8 posti, ed ha predisposto davanti ad ogni posto un
	cartellino con il nome di uno dei commensali.  Giunti al
	ristorante, gli 8 amici si siedono al tavolo a loro riservato, ma
	ancora in ordine casuale.
	
	Determinare la probabilit\`{a} che ognuno sieda al posto che reca 
	il suo nome, oppure in uno dei due posti adiacenti.
\end{enumerate}
}

\renewcommand{\Risposta}{
\begin{enumerate}
	\renewcommand{\labelenumi}{(\alph{enumi})} 
	\item  17/20160

	\item  7/5760
\end{enumerate}
}

\renewcommand{\Aiutino}{
\begin{enumerate}
	\renewcommand{\labelenumi}{(\alph{enumi})} 
	
	\item Generalizziamo il problema ad $n$ persone, ed indichiamo con
	$L_{n}$ il numero di permutazioni che verificano la propriet\`{a}
	richiesta.  Si tratta di trovare una relazione ricorrente per
	$L_{n}$.  Dove pu\`{o} sedersi il proprietario del biglietto
	numero~1?  Quante sono (in funzione dei casi precedenti) le
	disposizioni possibili se si siede al posto numero~1?
	
	Se invece il proprietario del biglietto numero~1 si siede al posto
	numero~2, chi si pu\`{o} sedere al posto numero~1?  Quante sono a
	quel punto le possibili permutazioni?
	
	\item Numeriamo opportunamente persone e posti.  L'idea \`{e} di
	ridursi al problema della fila considerando dove si siede il
	proprietario del posto numero~1.
	
	Se si siede al suo posto, allora ci siamo ridotti al problema su 
	un fila con una persona in meno (come? perch\'{e}?).
	
	Se si siede al posto numero~2, chi si pu\`{o} sedere al posto 
	numero~1? Ora ci sono due possibilit\`{a}: se \`{e} il numero~2, 
	allora ci siamo ridotti \ldots; se invece \`{e} il numero~8, 
	allora come devono sedersi gli altri?
	
	Quali altre possibilit\`{a} ci sono per il proprietario del posto 
	numero~1?
\end{enumerate}
}

\renewcommand{\SoluzioneUfficiale}{\TitoloSoluzione\DomandaA
La probabilit\`{a} richiesta \`{e} $34/8!=17/20160$.

I casi possibili sono tanti quante le possibili permutazioni di 8 
elementi, e cio\`{e} 8!.

Per contare i casi favorevoli, generalizziamo il problema ad un 
gruppo di $n$ persone, ed indichiamo con $L_{n}$ il numero di 
permutazioni con la propriet\`{a} richiesta. \`{E} banale 
verificare che $L_{1}=1$ ed $L_{2}=2$. Consideriamo ora un gruppo di 
$n$ persone, con $n\geq 3$. Il proprietario del biglietto numero~1 
pu\`{o} sedersi in 2 posti: il suo o quello adiacente.
\begin{itemize}
	\item  Se il proprietario del biglietto numero~1 si siede al suo 
	posto, ci siamo ridotti allo stesso problema con una persona in 
	meno, dunque abbiamo $L_{n-1}$ configurazioni possibili.

	\item  Se il proprietario del biglietto numero~1 si siede al 
	posto numero~2, allora necessariamente il possessore del 
	biglietto numero~2 deve sedersi al posto numero~1 (gli altri 
	infatti non possono farlo senza violare la condizione imposta 
	dal problema). Ne segue quindi che i possessori dei biglietti~1 
	e~2 si sono scambiati il posto, e ci siamo ridotti allo stesso 
	problema per gli $(n-2)$ possessori dei biglietti dal 3 in poi, 
	per i quali abbiamo quindi $L_{n-2}$ configurazioni possibili.
\end{itemize}

Abbiamo cos\`{i} dimostrato che 
$$L_{n}=L_{n-1}+L_{n-2}
\quad\quad
\forall n\geq 3.$$

Partendo da $L_{1}=1$ ed $L_{2}=2$ possiamo quindi facilmente 
ricavarci i valori successivi, che risultano essere $L_{3}=3$, 
$L_{4}=5$, $L_{5}=8$, $L_{6}=13$, $L_{7}=21$, $L_{8}=34$.

\DomandaB
La probabilit\`{a} richiesta \`{e} $49/8!=7/5760$.

I casi possibili sono ancora tanti quante le possibili permutazioni di 8 
elementi, e cio\`{e} 8!.

Per contare i casi favorevoli, indichiamo convenzionalmente gli amici 
ed i posti con i numeri da 1 ad 8, e distinguiamo i casi a seconda di 
dove siede il commensale numero~1.
\begin{itemize}
	\item  Se il commensale numero~1 si siede al proprio posto, 
	allora ci siamo ridotti al problema della fila sulle restanti 7 
	persone e posizioni, per cui abbiamo $L_{7}$ possibilit\`{a}.

	\item  Se il commensale numero~1 si siede al posto numero~2, 
	allora ci sono 2 sottocasi a seconda di chi occupa il posto 
	numero~1.
	\begin{itemize}
		\item Se il posto numero~1 viene occupato dal commensale
		numero~2, allora semplicemente i commensali~1 e~2 si sono
		scambiati il posto. Ci siamo quindi ridotti al problema della 
		fila sulle restanti 6 posizioni, per un totale di $L_{6}$ 
		possibilit\`{a}.
	
		\item Se il posto numero~1 viene occupato dal commensale
		numero~8, allora necessariamente il commensale~2 dovr\`{a}
		sedersi al posto~3 (non ha infatti altri posti a
		disposizione), il commensale~3 al posto~4, e cos\`{i} via.  In
		poche parole, la configurazione \`{e} fissa e si riduce ad una
		rotazione della tavolata.  In questo sottocaso abbiamo quindi
		una sola possibilit\`{a}.
	\end{itemize}

	\item  Se il commensale numero~1 si siede al posto numero~8, 
	allora la situazione \`{e} perfettamente simmetrica rispetto alla 
	precedente. Si generano dunque 2 sottocasi che producono, 
	rispettivamente, $L_{6}$ ed 1 configurazione. 
\end{itemize}

In conclusione, il numero di casi favorevoli \`{e} $L_{7}+2L_{6}+2=49$.
}

\renewcommand{\SoluzioniAlternative}{}

\renewcommand{\Osservazioni}{\TitoloOsservazione

Data la relazione ricorrente che soddisfa, la successione $L_{n}$ 
altro non \`{e} che la ben nota successione di Fibonacci.

Inoltre il ragionamento svolto al punto~(b) si generalizza facilmente
al caso di $n$ commensali, mostrando che in generale il numero di
disposizioni intorno al tavolo che rispettano la propriet\`{a} data
\`{e} $L_{n-1}+2L_{n-2}+2$.  }

\renewcommand{\Commenti}{}